% 翻译状态：专题标题已译；正文待继续翻译与校对。
% Chapter 3, Topic _Linear Algebra_ Jim Hefferon
%  http://joshua.smcvt.edu/linearalgebra
%  2001-Jun-12
\topic{线性映射的几何} 
\index{Geometry of Linear Maps|(}
%The geometry of linear maps %$\map{h}{\Re^n}{\Re^m}$ 
%is appealing both for its simplicity 和 for its usefulness.
下面两组图对比了非线性映射
\( f_1(x)=e^x \) 和 \( f_2(x)=x^2 \) 
\begin{center}
  \includegraphics{map/mp/ch3.46}
  \hspace*{4em}
  \includegraphics{map/mp/ch3.47}
\end{center}
与线性映射
\( h_1(x)=2x \) 和 \( h_2(x)=-x \).
\begin{center}
  \includegraphics{map/mp/ch3.48}
  \hspace*{4em}
  \includegraphics{map/mp/ch3.49}
\end{center}
四幅图都把左侧的定义域 $\Re$
映到右侧的陪域 $\Re$。 
箭头标出每个映射对
$x=0$, $x=1$, $x=2$, $x=-1$, 和 $x=-2$.

非线性映射将定义域变换为值域时会扭曲定义域。
例如,
\( f_1(1) \) 与
$f_1(2)$ 的距离大于它与 $f_1(0)$ 的距离 \Dash 该映射使
定义域不均匀地扩张，使得 $x=2$ 附近的区间
扩张得更大，
而 $x=0$ 附近的区间扩张得更小。
线性映射则更规则：
因为对每个映射，整个定义域
都按相同因子扩张。
左侧的映射 $h_1$ 使所有区间扩张为原来的两倍，
而右侧的 $h_2$ 保持区间长度不变但反转
它们的方向，例如将从 $1$ 到 $2$ 的上升区间
变换为
从 $-1$ 到~$-2$ 的下降区间。

从 $\Re$ 到 $\Re$ 的线性映射只有标量乘法，但
在更高维空间中还会发生更多情况。
例如，这个 $\Re^2$ 上的线性变换使向量逆时针旋转。\index{rotation}\index{linear map!rotation}
\begin{center}
  \includegraphics{map/mp/ch3.50}
\end{center}
将向量投影到 $xz$ 平面的 $\Re^3$ 上的变换
将向量投影到 $xz$ 平面
也不只是简单的缩放。
\begin{center}
 \includegraphics{map/mp/ch3.51}
\end{center}

尽管形式更加多样，
即使在高维空间中，线性映射的行为也很好。
考虑线性映射 $\map{h}{\Re^n}{\Re^m}$ 和
用标准基将其表示为矩阵 $H$。
回忆定理 V.\ref{th:CanonFormForMatEquiv}：
$H$ 可分解为 $H=PBQ$，
其中 $P,Q$ 非奇异，而 $B$ 是部分恒等矩阵。
还要回忆非奇异矩阵
可分解为初等
矩阵\index{matrix!elementary reduction}\index{elementary reduction matrix}
$PBQ=T_nT_{n-1}\cdots T_sBT_{s-1}\cdots T_1$,
这些矩阵
由单位矩阵 $I$ 经过一次高斯行变换得到，
因此每个 $T$ 矩阵属于以下三种之一
\begin{equation*}
  I\grstep{k\rho_i}M_i(k) 
  \qquad 
  I\grstep{\rho_i\leftrightarrow\rho_j}P_{i,j}  
  \qquad
  I\grstep{k\rho_i+\rho_j}C_{i,j}(k) 
\end{equation*}
其中 $i\neq j$、$k\neq0$。
因此，若理解由部分恒等矩阵描述的线性映射
以及由初等矩阵描述的线性映射
的几何效果，就可以在某种意义上
完全理解任意线性映射的效果。
（为便于绘图，下面的图只展示 $\Re^2$ 上的变换
但这些原理适用于从任意 $\Re^n$ 到任意 $\Re^m$ 的映射。）

由部分恒等矩阵表示的线性变换的几何效果是投影。

\begin{equation*}
  \colvec{x \\ y  \\ z}
  \quad\xrightarrow{{\text{\tiny$\displaystyle\left(\begin{array}[r]{@{}c@{\hspace{0.8em}}c@{\hspace{0.8em}}c@{}}
                1  &0  &0 \\
                0  &1  &0 \\
                0  &0  &0
          \end{array}\right)$}}}
% \begin{mat}[r]
%       1  &0  &0   \\
%       0  &1  &0   \\
%       0  &0  &0   
%     \end{mat}} % _{\stdbasis_3,\stdbasis_3}}
  \qquad
  \colvec{x \\ y  \\ 0}
\end{equation*}

矩阵 $M_i(k)$ 的几何效果是

沿第 $i$ 个坐标轴将向量按因子 $k$ 拉伸。
该映射沿 $x$ 轴按因子 $3$ 拉伸。
\begin{center}
  \includegraphics{map/mp/ch3.52}
\end{center}
若 $0\leq k<1$ 或 $k<0$，则第 $i$ 个
分量反向变化，此处向左。
\begin{center}
  \includegraphics{map/mp/ch3.53}
\end{center}
上述任一种拉伸都称为
\definend{伸缩}.\index{dilation}\index{linear map!dilation}

由矩阵 $P_{i,j}$ 表示的变换
交换第 $i$ 个和第 $j$ 个坐标轴。
这是 \definend{反射}\index{reflection}\index{linear map!reflection} 
关于直线 $x_i=x_j$ 的反射。
\begin{center}
  \includegraphics{map/mp/ch3.54}
\end{center}
涉及两个以上坐标轴的置换可分解为多个坐标轴对交换的组合；
坐标轴对交换的组合；见习题 \nearbyexercise{exer:PermIsCompSwaps}。

其余矩阵都具有 $C_{i,j}(k)$ 的形式。
例如，$C_{1,2}(2)$ 执行 $2\rho_1+\rho_2$。 
\begin{equation*}
  \colvec{x  \\  y}\quad
  \xrightarrow{\text{\tiny$\displaystyle\left(\begin{array}[r]{@{}c@{\hspace{0.8em}}c@{}}
                1  &0  \\
                2  &1  
          \end{array}\right)$}}  % _{\stdbasis_2,\stdbasis_2}}
   \quad\colvec{x \\ 2x+y}
\end{equation*}
在下图中，
第一分量为 $1$ 的向量 $\vec u$ 受到的影响小于
第一分量为 $2$ 的向量 $\vec v$。
向量 $\vec u$ 映到
仅比 $\vec u$ 高 $2$ 的 $h(\vec u)$，而
$h(\vec v)$ 比 $\vec v$ 高 $4$。
\begin{center}
  \includegraphics{map/mp/ch3.55}
\end{center}
任何第一分量为 $1$ 的向量都会以
与 $\vec u$ 相同的方式受影响：
它向上滑动 $2$。
任何第一分量为 $2$ 的向量都会向上滑动 $4$，
正如 $\vec v$ 一样。
也就是说，由
$C_{i,j}(k)$ 表示的变换根据向量的第 $i$ 个分量影响向量。

还可以考察该映射
在单位正方形上的作用。
在下一幅图中，
第一分量为 $0$ 的向量（如原点）不会被
竖直推动，而第一分量为正的向量会向上滑动。
这里，第一分量为 $1$ 的所有向量，即正方形的整条
右边，都滑动相同的距离。
一般地，位于同一条竖直线上的向量平移量相同，
滑动量是其第一分量的两倍。
所得图形与正方形具有相同的底和高（因而面积相同），但直角顶点消失了。
\begin{center}
  \includegraphics{map/mp/ch3.56}
\end{center}

作为对比，下一幅图展示矩阵
$C_{2,1}(2)$.
表示的映射效果。这里向量根据其
第二分量受到影响：
$\binom{x}{y}$沿水平方向滑动 $2y$。
\begin{center}
  \includegraphics{map/mp/ch3.57}
\end{center}
一般地，对于任意 $C_{i,j}(k)$，第 $i$ 个分量相同的向量平移量相同。
这种映射称为
\definend{剪切}.\index{shear}\index{linear map!shear}

这样我们就理解了右侧四类矩阵的几何效果，
右侧矩阵中的
$H=T_nT_{n-1}\cdots T_jBT_{j-1}\cdots T_1$
因而在某种意义上理解了
任意矩阵 $H$ 的作用。
因此，即使在高维空间中，线性映射的几何也很简单：它
由若干个
以简单方式作用的分量组合而成。

我们将从两方面应用这种理解。
第一方面是证明线性映射几何的一个一般结论。
线性映射的几何。
回忆在线性映射下，子空间的像仍是子空间，
因此由 $H$ 表示的线性变换 $h$ 将经过原点的直线
映到经过原点的直线。
（像空间的维数不可能大于
定义域空间的维数，所以直线不可能映成平面。）
我们将证明 $h$ 把任意直线 \Dash 不仅是经过原点的直线 \Dash 
映成直线。
证明很简单：
部分恒等投影 $B$ 和初等矩阵 $T_i$
都把直线输入变成直线输出；
四种情形的验证见习题 \nearbyexercise{exer:ImageLinSurIsLinSur}。
因此它们的复合也保持直线。
% Thus, by underst和ing its components we can underst和 arbitrary square 
% matrices $H$, in the sense that we can prove things about them.

我们应用这种理解的第二个方面是
线性映射几何的理解，用来
阐明微积分中的一个要点。
下面的图展示
单变量实函数 \( y(x)=x^2+x \) 的作用。
与前面所示的非线性函数一样，
该映射的几何效果
并不规则，因为在定义域的不同点作用不同；例如，
例如，当输入 $x$ 从 $2$ 变为 $-2$ 时，相应输出 $f(x)$
先减小，短暂停留，
然后增加。
\begin{center}
  \includegraphics{map/mp/ch3.58}
\end{center}
但在微积分中，我们较少关注映射整体，而更多关注
映射的局部作用。
下面仔细观察该映射
在 $x=1$ 附近的作用。
导数为 $dy/dx=2x+1$，
因此在 \( x=1 \) 附近
有 \( \Delta y\approx 3\cdot\Delta x \)。
%; in other words, \( (1.001^2+1.001)-(1^2+1)\approx 3\cdot (0.001) \).
也就是说，在 $x=1$ 的邻域内，
定义域经该映射变换后扩张了
因子 $3$ \Dash  局部而言，
近似地说，这是一次伸缩。
%The map $y$ is locally regular.
下图将其表示为定义域中的小区间
定义域区间 $(1-\Delta x\,..\,1+\Delta x)$
映到陪域中的区间 $(2-\Delta y\,..\,2+\Delta y)$
其宽度约为前者的三倍。 % $\Delta y \approx 3\cdot \Delta x$.
\begin{center}
  \includegraphics{map/mp/ch3.59}
\end{center}
% (When the above picture is drawn in the traditional cartesian way
% then the prior sentence about the rate of growth of $y(x)$ is usually
% stated:~the derivative $3$ is the slope of the 
% line tangent to the graph at the point $(1,2)$.)

在高维空间中，核心思想相同，但可能出现更多情况。
对于函数 \( \map{y}{\Re^n}{\Re^m} \) 和点 \( \vec{x}\in\Re^n \),
导数定义为
线性映射 \( \map{h}{\Re^n}{\Re^m} \)，最佳近似 $y$ 在 $y(\vec{x})$ 附近的变化。
how \( y \) changes near \( y(\vec{x}) \).
因此，上述几何描述可直接应用于导数。
%Calculus considers the map 
%that locally approximates the change \( \Delta x\mapsto 3\cdot\Delta x \)
%(instead of the actual change map
%\( \Delta x\mapsto y(1+\Delta x)-y(1) \)) because the local map 
%is easier.
%It is easier in that, if the
%input change is doubled, or tripled, etc., then the
%resulting output change will double, or triple, etc.
%\begin{equation*}
%  3(r\,\Delta x)=r\,(3\Delta x)
%\end{equation*}
%(for $r\in\Re$), 
%和 it is easier in that 
%adding changes in input adds the resulting output changes.
%\begin{equation*}
%  3(\Delta x_1+\Delta x_2)=3\Delta x_1+3\Delta x_2
%\end{equation*}
%In short, what's
%easy about \( \Delta x\mapsto 3\cdot\Delta x \) is that
%it 是线性的.

最后说明这一观点如何
澄清一个常被误解的
关于导数的结果：链式法则。
回忆在两个函数满足适当条件时，
复合函数的导数为：
\begin{equation*}
  \frac{d\,(\composed{g}{f})}{dx}(x) = 
  \frac{dg}{dx}(f(x))\cdot\frac{df}{dx}(x)
\end{equation*} 
例如，$\sin(x^2+3x)$ 的导数为
$\cos(x^2+3x)\cdot(2x+3)$.

这从何而来？
考虑 $\map{f,g}{\Re}{\Re}$。 
\begin{center}
  \includegraphics{map/mp/ch3.60}
\end{center}
第一个映射 $f$ 将 $x$ 的邻域伸缩因子设为
\begin{equation*}
  \frac{df}{dx}(x) 
\end{equation*}
第二个映射 $g$ 随后将 
将 $f(x)$ 的邻域伸缩因子设为
\begin{equation*}
  \frac{dg}{dx}(\,f(x)\,) 
\end{equation*}
二者复合后， 
复合映射的伸缩因子就是二者的乘积。
在高维空间中，
描述函数在某点附近变化的映射是线性映射，
并由矩阵表示。
链式法则就是将这些矩阵相乘。

% Thus, the geometry of linear maps $\map{h}{\Re^n}{\Re^m}$ 
% is appealing both for its simplicity 和 for its usefulness.

\begin{exercises}
  \item 
    Use the $H=PBQ$ decomposition to find
    伸缩、翻转、剪切和投影的组合
    生成映射 $\map{h}{\Re^3}{\Re^3}$
    相对于标准基由下列矩阵表示。
    \begin{equation*}
      H=\begin{mat}[r]
        1  &2  &1  \\
        3  &6  &0  \\
        1  &2  &2
      \end{mat}
    \end{equation*}
    \begin{answer}
      下面的高斯消元
      \begin{multline*}
        \grstep[-\rho_1+\rho_3]{-3\rho_1+\rho_2}
        \begin{mat}[r]
          1  &2  &1  \\
          0  &0  &-3 \\
          0  &0  &1
        \end{mat}
        \grstep{(1/3)\rho_2+\rho_3}
        \begin{mat}[r]
          1  &2  &1  \\
          0  &0  &-3 \\
          0  &0  &0
        \end{mat}                                     \\
        \grstep{(-1/3)\rho_2}
        \begin{mat}[r]
          1  &2  &1  \\
          0  &0  &1 \\
          0  &0  &0
        \end{mat}
        \grstep{-\rho_2+\rho_1}
        \begin{mat}[r]
          1  &2  &0  \\
          0  &0  &1 \\
          0  &0  &0
        \end{mat}
      \end{multline*}
      得到该矩阵的简化阶梯形。
      接着，对列作两次操作：先将第一列乘以 $-2$
      加到第二列，再交换第二、三列
      得到如下部分恒等矩阵。
      \begin{equation*} 
        B=\begin{mat}[r]
          1  &0  &0  \\
          0  &1  &0  \\ 
          0  &0  &0
        \end{mat}
      \end{equation*}
      以上内容用矩阵表示为：
      \begin{equation*}
        P=
        \begin{mat}[r]
          1  &-1    &0  \\
          0  &1     &0  \\
          0  &0     &1         
        \end{mat}
        \begin{mat}[r]
          1  &0    &0  \\
          0  &-1/3 &0  \\
          0  &0    &1
        \end{mat}
        \begin{mat}[r]
          1  &0    &0  \\
          0  &1    &0  \\
          0  &1/3  &1         
        \end{mat}
        \begin{mat}[r]
          1  &0  &0  \\
          0  &1  &0  \\
         -1  &0  &1         
        \end{mat}
        \begin{mat}[r]
          1  &0  &0  \\
         -3  &1  &0  \\
          0  &0  &1         
        \end{mat}
      \end{equation*}
      和 
      \begin{equation*}
        Q=
        \begin{mat}[r]
          1  &-2    &0  \\
          0  &1     &0  \\
          0  &0     &1         
        \end{mat}
        \begin{mat}[r]
          0  &1     &0  \\
          1  &0     &0  \\
          0  &0     &1         
        \end{mat}
      \end{equation*}
      因此给定矩阵可分解为 $PBQ$。
    \end{answer}
  \item 
    哪些伸缩、翻转、剪切和投影的组合
    能产生逆时针旋转 $2\pi/3$ 弧度？
    \begin{answer}
      首先用矩阵 $H$ 表示该映射，
      进行行变换，并在需要时进行列变换
      将其化为部分恒等矩阵。
      然后将其写成分解 $H=PBQ$。
      代入旋转矩阵的一般形式
          \begin{equation*}
            \rep{r_\theta}{\stdbasis_2,\stdbasis_2}
            \begin{mat}
              \cos\theta  &-\sin\theta  \\
              \sin\theta  &\cos\theta
            \end{mat}
          \end{equation*}
          得到如下表示。
          \begin{equation*}
            \rep{r_{2\pi/3}}{\stdbasis_2,\stdbasis_2}
            \begin{mat}[r]
              -1/2        &-\sqrt{3}/2  \\
              \sqrt{3}/2  &-1/2
            \end{mat}
          \end{equation*}
          Gauss's Method is routine.
          \begin{equation*}
            \grstep{\sqrt{3}\rho_1+\rho_2}
            \begin{mat}[r]
              -1/2        &-\sqrt{3}/2  \\
               0          &-2
            \end{mat}
            \grstep[(-1/2)\rho_2]{-2\rho_1}
            \begin{mat}[r]
               1          &\sqrt{3}    \\
               0          &1
            \end{mat}
            \grstep{-\sqrt{3}\rho_2+\rho_1}
            \begin{mat}[r]
               1          &0   \\
               0          &1
            \end{mat}
          \end{equation*}
          这对应于如下矩阵方程。
          \begin{equation*}
            \begin{mat}[r]
              1  &-\sqrt{3}  \\
              0  &1
            \end{mat}
            \begin{mat}[r]
              -2  &0    \\
               0  &-1/2
            \end{mat}
            \begin{mat}[r]
               1         &0  \\
               \sqrt{3}  &1
            \end{mat}
            \begin{mat}[r]
              -1/2        &-\sqrt{3}/2  \\
              \sqrt{3}/2  &-1/2
            \end{mat}
            =I
          \end{equation*}
          对方程取逆并解出 $H$，得到以下分解。
          \begin{equation*}
            \begin{mat}[r]
              -1/2        &-\sqrt{3}/2  \\
              \sqrt{3}/2  &-1/2
            \end{mat}
            =
            \begin{mat}[r]
                1         &0  \\
               -\sqrt{3}  &1
            \end{mat}
            \begin{mat}[r]
              -1/2  &0    \\
               0    &-2
            \end{mat}
            \begin{mat}[r]
              1  &\sqrt{3}  \\
              0  &1
            \end{mat}
            I
          \end{equation*}
    \end{answer}
  \item 
    若映射非奇异，从其表示化为
    单位矩阵无需任何列变换，
    因此在 $H=PBQ$ 中矩阵 $Q$ 是单位矩阵。
    非奇异映射的一个例子是
    变换 $\map{t_{-\pi/4}}{\Re^2}{\Re^2}$，它将
    向量顺时针旋转 $\pi/4$ 弧度。
    \begin{exparts}
      \partsitem 求出表示
         相对于标准基表示该映射。
      \partsitem 用高斯-约旦法将 $H$ 化为单位矩阵，
        不使用列变换将 $H$ 化为单位矩阵。
      \partsitem
        将其写成矩阵方程
        $T_jT_{j-1}\cdots T_1H=I$.
      \partsitem 解矩阵方程得到 $H$。
      \partsitem 将 $H$ 描述为
        伸缩、翻转、剪切和投影的组合
        （单位矩阵是平凡投影）。
    \end{exparts}
    \begin{answer}
      \begin{exparts}
        \partsitem 回忆，逆时针旋转
          $\theta$ 弧度相对于标准基表示为
          可表示为：
          \begin{equation*}
            \rep{t_{\pi/4}}{\stdbasis_2,\stdbasis_2}
            =\begin{mat}
              \cos\theta  &-\sin\theta  \\
              \sin\theta  &\cos\theta
             \end{mat}
          \end{equation*}
          顺时针角是逆时针角的相反数，
          one.  
          \begin{equation*}
            \rep{t_{-pi/4}}{\stdbasis_2,\stdbasis_2}
            =\begin{mat}
              \cos(-\pi/4)  &-\sin(-\pi/4)  \\
              \sin(-\pi/4)  &\cos(-\pi/4)
            \end{mat}
            =\begin{mat}[r]
              \sqrt{2}/2  &\sqrt{2}/2  \\
              -\sqrt{2}/2 &\sqrt{2}/2
            \end{mat}
          \end{equation*}
        \partsitem
          下面的高斯-约旦消元
          \begin{equation*}
            \grstep{\rho_1+\rho_2}
            \begin{mat}[r]
              \sqrt{2}/2  &\sqrt{2}/2  \\
              0           &\sqrt{2}
            \end{mat}
            \grstep[(1/\sqrt{2})\rho_2]{(2/\sqrt{2})\rho_1}
            \begin{mat}[r]
              1  &1  \\
              0  &1
            \end{mat}
            \grstep{-\rho_2+\rho_1}
            \begin{mat}[r]
              1  &0  \\
              0  &1
            \end{mat}
          \end{equation*}
          得到单位矩阵。
          因此无需交换列
          即可得到部分恒等矩阵。
        \partsitem In matrix multiplication the reduction is
          \begin{equation*}
            \begin{mat}[r]
              1  &-1 \\
              0  &1
            \end{mat}
            \begin{mat}[r]
              2/\sqrt{2}  &0         \\
              0           &1/\sqrt{2}
            \end{mat}
            \begin{mat}[r]
              1  &0 \\
              1  &1
            \end{mat}
            H
            =I
          \end{equation*}
          （注意，高斯操作的复合从右向左进行）。
        \partsitem  Taking inverses 
          \begin{equation*}
            H
            =
            \underbrace{
              \begin{mat}[r]
                1  &0 \\
                -1  &1
              \end{mat}
              \begin{mat}[r]
                \sqrt{2}/2  &0         \\
                0           &\sqrt{2}
              \end{mat}
              \begin{mat}[r]
                1  &1 \\
                0  &1
              \end{mat}
             }_P
            I
          \end{equation*}
          得到 $H$ 所需的分解。
          部分恒等矩阵就是 $I$。
          % , 和 $Q$ is trivial, that is, it is also an identity
          % matrix.
        \partsitem 从右向左读取复合（忽略
          单位矩阵视为平凡矩阵），可知 $H$ 的作用等同于
          先执行以下剪切
          \begin{center}
            \includegraphics{map/mp/ch3.95}
         \end{center}
         再进行伸缩，将所有第一分量乘以
         $\sqrt{2}/2$（这是收缩，因为 $\sqrt{2}/2\approx0.707$ 
         小于 $1$）。
         并将所有第二分量乘以 $\sqrt{2}$，
         最后再进行一次剪切。
          \begin{center}
            \includegraphics{map/mp/ch3.96}
         \end{center}
         例如，取与
         $x$ 轴夹角为 $\pi/6$ 的单位向量，依次应用 $H$ 的各个分量。
          \begin{center}
            \includegraphics{map/mp/ch3.97}
         \end{center}
         容易验证所得向量长度为 1，
         与 $x$ 轴夹角为 $-\pi/12$，确实形成
         即顺时针旋转 $\pi/4$ 弧度，
         因为 $(\pi/6)-(\pi/4)=-\pi/12$。 
      \end{exparts}
    \end{answer}
  \item \label{exer:RToRIsScalMult} 
    证明 $\Re^1$ 上任意线性变换都是某个映射 $h_k$
    将标量乘法 $x\mapsto kx$。
    \begin{answer}
      相对于
      标准基 $\stdbasis_1,\stdbasis_1$。 
      这得到一个 $\nbyn{1}$ 矩阵。
      唯一元素就是标量~$k$。
    \end{answer}
  \item \label{exer:ImageLinSurIsLinSur} 
    证明线性映射保持空间的线性结构。
    \begin{exparts}
      \partsitem 证明从 $\Re^n$ 到 $\Re^m$ 的任意线性映射，
         任一直线的像仍是直线。
         该像可能是退化直线，即单个点。
      \partsitem 证明任意线性曲面的像仍是线性曲面。
         这推广了线性映射下
         子空间的像仍是子空间这一结果。
      \partsitem 线性映射还保持其他线性性质。
         证明线性映射保持“介于其间”关系：若点
         若 $B$ 位于 $A$ 与 $C$ 之间，则 $B$ 的像位于
         $A$ 的像与 $C$ 的像之间。
    \end{exparts}
    \begin{answer}
      \begin{exparts}
        \partsitem 直线是 $\Re^n$ 中形如
          $\set{\vec{v}=\vec{u}+t\cdot\vec{w}\suchthat t\in\Re}$.
          该直线上点的像为
          $h(\vec{v})=h(\vec{u}+t\cdot\vec{w})=h(\vec{u})+t\cdot h(\vec{w})$,
          当 $t$ 遍历实数时，这些向量的集合是
          一条直线（若 $h(\vec{w})=\zero$ 则为退化直线）。
        \partsitem 这是前面论证的直接推广。
        \partsitem 若点~$B$ 位于点~$A$ 与~$C$ 之间，则
          从 $A$ 到 $C$ 的线段包含 $B$。
          也就是说，存在 $t\in (0\,..\,1)$ 使得
          $\vec{b}=\vec{a}+t\cdot (\vec{c}-\vec{a})$（其中 $B$ 是
          （其中 $B$ 是 $\vec b$ 的端点，其他类似）。
          如第一小题的论证，线性性说明
          $h(\vec{b})=h(\vec{a})+t\cdot h(\vec{c}-\vec{a})$.  
      \end{exparts}
    \end{answer}
  \item 
    参考链式法则讨论中的图回答：
    若函数 $\map{f}{\Re}{\Re}$ 可逆，函数局部近似伸缩空间的倍数与其逆映射伸缩空间的倍数有何关系？
    \begin{answer}
      二者互为倒数。
      例如，对固定的 $x\in\Re$，
      若 $f^\prime (x)=k$（其中 $k\neq 0$），则 
      $(f^{-1})^\prime (x)=1/k$.
      \begin{center}
        \includegraphics{map/mp/ch3.98}
     \end{center}
    \end{answer}
  \item \label{exer:PermIsCompSwaps}
    证明这些数的任意排列（即重新排序）$p$
    $1$, \ldots, $n$ 的排列诱导的映射 
    \begin{equation*}
      \colvec{x_1 \\ x_2 \\ \vdots\\ x_n}
      \mapsto
      \colvec{x_{p(1)} \\ x_{p(2)} \\ \vdots \\ x_{p(n)}}
    \end{equation*}
    可由映射复合完成，
    其中每个映射只交换一对坐标。
    \textit{Hint:} you can use induction on $n$.
    （\textit{注：}第四章将证明这一点，并且还将
    证明交换次数的奇偶性由 $p$ 决定。
    也就是说，尽管某个
    排列可用两种不同方式表示，
    交换次数也不同，但两种方式要么都使用偶数次
    交换，要么都使用奇数次。）
    \begin{answer}
      对向量分量数作归纳即可证明。
      vector.
      In the $n=1$ base case the only permutation is the trivial one,
      且映射
      \begin{equation*}
        \colvec{x_1}
        \mapsto
        \colvec{x_1}
      \end{equation*}
      可表示为交换的复合，即零次交换。
      归纳步骤中，假设由
      少于 $n$ 个数的任意排列诱导的映射
      $n$ 个数的排列可仅用交换表示；现在考虑映射
      由
      $n$ 个数的排列 $p$ 诱导。
      \begin{equation*}
        \colvec{x_1 \\ x_2 \\ \vdots \\ x_n}
        \mapsto
        \colvec{x_{p(1)} \\ x_{p(2)} \\ \vdots \\ x_{p(n)}}
      \end{equation*}
      考虑满足 $p(i)=n$ 的编号 $i$。
      映射
      \begin{equation*}
        \colvec{x_1      \\ x_2      \\ \vdots \\ x_i      \\ \vdots \\ x_n}
        \mapsunder{\hat{p}}
        \colvec{x_{p(1)} \\ x_{p(2)} \\ \vdots \\ x_{p(n)} \\ \vdots  \\ x_{n}}
      \end{equation*}
      再交换第 $i$ 个与第 $n$ 个分量后，
      得到映射~$p$。
      根据归纳假设，$\hat{p}$ 可表示为交换的复合。
      的交换复合。
    \end{answer}
\end{exercises}
\index{Geometry of Linear Maps|)}
\endinput
